\documentclass[11pt]{article}
\usepackage[margin=1in]{geometry}
\usepackage{amsmath,amssymb}
\usepackage{booktabs}
\usepackage{array}
\usepackage{graphicx}
\usepackage{hyperref}
\usepackage[table]{xcolor}
\usepackage{microtype}
\usepackage{caption}

\hypersetup{
  colorlinks=true,
  linkcolor=blue!60!black,
  urlcolor=blue!60!black
}

\newcommand{\win}{\cellcolor{green!15}}
\newcommand{\lose}{\cellcolor{red!15}}
\newcommand{\tie}{\cellcolor{gray!10}}

\title{Random Augmentation vs.\ Canonicalization\\for Laplacian Positional Encodings\\
       \large Does test-time ambiguity-group averaging help?}
\author{ogbg-molpcba experiments, GIN backbone}
\date{\today}

\begin{document}
\maketitle

\begin{abstract}
We compare three ways of resolving the sign + basis ambiguity of Laplacian
eigenvector positional encodings (LapPE) on the ogbg-molpcba molecular
benchmark: (i) deterministic canonicalization (\texttt{map}, \texttt{oap}),
(ii) training-time random augmentation (\texttt{random\_augmented}), and
(iii) random augmentation \emph{plus} test-time averaging over $K$ random
draws of the eigenvector ambiguity group. Two sweeps (54 runs total) at
$k\!\in\!\{3,8\}$ eigenvectors, hidden dim
$h\!\in\!\{16,128,512\}$, three seeds each. \textbf{Headline}: at $k\!=\!3$
canonicalization (\texttt{oap}) wins at $h\!=\!512$; at $k\!=\!8$ random
augmentation wins at \emph{every} hidden dim. Test-time averaging adds a
tiny consistent bump ($+0.0001$ to $+0.0005$ AP), never hurts.
\textbf{The single-line conclusion: canonicalization does \emph{not} perform
better, except in a narrow $k\!=\!3$ / $h\!=\!512$ corner.}
\end{abstract}

\tableofcontents
\bigskip

%==============================================================================
\section{What problem are we solving?}
%==============================================================================

Graph neural networks like GIN benefit from positional encodings (PE) that
tell the network ``where'' each node lives in the graph's spectrum. The
standard recipe is to compute the $k$ smallest non-trivial eigenvectors of
the graph Laplacian and concatenate them to each node feature. This works,
but eigenvectors are not unique: for any simple eigenvalue, the eigenvector
is defined only up to a sign $(\pm 1)$; for a multiplicity-$m$ eigenvalue,
\emph{any orthogonal change of basis $Q\in O(m)$} produces an equally valid
eigenbasis. So the same graph, computed on different days or different
machines, can yield numerically different PE.

Three ways to handle this ambiguity:

\paragraph{Canonicalization.} Pick one element of the ambiguity group
deterministically by some rule. Examples:
\begin{itemize}\itemsep0pt
  \item \texttt{map} --- maximum-absolute-position rule per column.
  \item \texttt{oap} --- ordered-absolute-position rule, more robust at
        multiplicity blocks.
\end{itemize}
The model sees the \emph{same} PE every epoch and at test time.

\paragraph{Random augmentation.}
(\texttt{random\_augmented}.) Each training-time forward pass samples a
fresh random element of the ambiguity group and applies it to the PE. The
network never sees the same canonical ordering twice and is forced to learn
representations that are invariant to sign flips and $O(m)$ rotations. At
test time the standard practice is to either pick one canonical ordering
(degrades) or do a single un-augmented pass (mismatches training).

\paragraph{Random augmentation $+$ test-time averaging.}
(\texttt{random\_augmented} with $K$-sample averaging at test time.) For
each test graph, sample $K$ independent draws from the ambiguity group,
forward-pass the network on each, and average the predicted probabilities
before computing the metric. This is the test-time analog of training-time
augmentation --- you take an empirical Monte-Carlo estimate of the
\emph{group-invariant} prediction. Inspired by the
\texttt{test\_with\_randomized\_invariance} helper in
\href{https://github.com/jwsiegel2510/Sn-invariant-weighted-frames}%
{Siegel \emph{et al.}, weighted-frames}; we adapt it from $S_n$ to the
LapPE ambiguity group $\bigtimes_i \mathrm{O}(m_i)$.

%==============================================================================
\section{Exact parameter settings}
%==============================================================================

Both sweeps share everything except $k$ and the test-time sample budget $K$.

\begin{table}[h]
\centering\small
\renewcommand{\arraystretch}{1.05}
\begin{tabular}{l l}
\toprule
\textbf{Knob}                       & \textbf{Value}                                                  \\
\midrule
Dataset                             & \texttt{ogbg-molpcba} (molecular property prediction)             \\
Backbone                            & GIN, 5 layers                                                    \\
PE type                             & eigvec (Laplacian eigenvectors)                                  \\
Eigval scaling                      & off                                                              \\
Hidden dim $h$                      & $\{16,\ 128,\ 512\}$                                              \\
Number of eigenvectors $k$          & $3$ (sweep~1) ~/~ $8$ (sweep~2)                                  \\
Cache eigenvectors                  & $15$ (slice to $k$ at runtime)                                    \\
Batch size                          & $32$                                                              \\
Learning rate                       & $1\!\times\!10^{-3}$                                              \\
Optimizer / scheduler               & Adam, default cosine                                              \\
Epoch budget (cap)                  & $200$                                                             \\
Early stopping patience             & $15$ (validation AP, stop after no-improve)                       \\
Seeds                               & $\{0,1,2\}$ ($3$ seeds per cell)                                  \\
Canonicalizations                   & \texttt{map}, \texttt{oap}, \texttt{random\_augmented}            \\
Test-time samples $K$               & $8$ (sweep~1) ~/~ $16$ (sweep~2)                                  \\
\midrule
Total runs                          & $3 \times 3 \times 3 = 27$ per sweep, $54$ total                  \\
Hardware                            & 2 $\times$ NVIDIA A40 (48~GB), $6$ workers per GPU                \\
\bottomrule
\end{tabular}
\caption{Hyperparameters held constant between the two sweeps. Bumping
$k$ from $3$ to $8$ slightly increases parameter count: at $h\!=\!512$,
$2{,}973{,}312$ vs $2{,}975{,}872$ params. The test-time sample budget $K$
is set to the size of the (simple-spectrum) sign group at $k\!=\!3$ and
to a $\sim$6.25\,\% Monte-Carlo coverage at $k\!=\!8$ where the group is
$\mathbb{Z}_2^8 = 256$.}
\label{tab:hparams}
\end{table}

\paragraph{Reproducibility.}
The augmentation rng is seeded as
$(\text{base\_seed} \cdot 1{,}000{,}003 + s) \times 2{,}654{,}435{,}761
\oplus \text{graph\_idx}$, so the aug-averaged metric is bit-reproducible
for a given $(\text{seed}, K)$.

%==============================================================================
\section{Results}
%==============================================================================

\subsection{Sweep~1: $k=3$, $K=8$}

At $k\!=\!3$ all three eigenvalues are simple (multiplicity~$1$) on
ogbg-molpcba molecules, so the ambiguity group is just
$\mathbb{Z}_2^3 = 8$. We picked $K\!=\!8$ to cover it exactly in
expectation.

\begin{table}[h]
\centering\small
\begin{tabular}{l r r r r r}
\toprule
canonicalization     &  $h$  & params     & AP (single pass)     & AP (aug-avg $K{=}8$) & $\Delta$ aug \\
\midrule
\tie map             &   16  &      5,744 & 0.0626 $\pm$ 0.0025  & ---                  & ---           \\
\tie oap             &   16  &      5,744 & 0.0641 $\pm$ 0.0051  & ---                  & ---           \\
\win random\_augmented &   16  &      5,744 & \textbf{0.0662} $\pm$ 0.0015 & 0.0663 $\pm$ 0.0015 & $+$0.0001    \\
\midrule
\tie map             &  128  &    202,752 & 0.1704 $\pm$ 0.0060  & ---                  & ---           \\
\tie oap             &  128  &    202,752 & 0.1682 $\pm$ 0.0089  & ---                  & ---           \\
\win random\_augmented &  128  &    202,752 & \textbf{0.1762} $\pm$ 0.0064 & 0.1765 $\pm$ 0.0062 & $+$0.0002    \\
\midrule
\tie map             &  512  &  2,973,312 & 0.2397 $\pm$ 0.0122  & ---                  & ---           \\
\win oap             &  512  &  2,973,312 & \textbf{0.2431} $\pm$ 0.0022 & ---          & ---           \\
\lose random\_augmented &  512  &  2,973,312 & 0.2395 $\pm$ 0.0117  & 0.2398 $\pm$ 0.0116 & $+$0.0002    \\
\bottomrule
\end{tabular}
\caption{Sweep~1, $k=3$, ogbg-molpcba, mean $\pm$ std over 3 seeds. Bold
numbers are the best mean AP at each $h$. \texttt{random\_augmented} wins
at $h\!=\!16$ and $h\!=\!128$; \texttt{oap} wins at $h\!=\!512$ by
$\mathbf{+0.0033}$ AP. Test-time averaging $K\!=\!8$ adds at most
$+0.0002$ AP --- it does not close the gap to \texttt{oap}.}
\label{tab:k3}
\end{table}

\subsection{Sweep~2: $k=8$, $K=16$}

At $k\!=\!8$ the sign group is $\mathbb{Z}_2^8 = 256$. We picked
$K\!=\!16$ to keep test-eval cost reasonable (Monte-Carlo error
$\propto 1/\sqrt{K} = 1/4$) while still being a meaningful sample.

\begin{table}[h]
\centering\small
\begin{tabular}{l r r r r r}
\toprule
canonicalization     &  $h$  & params     & AP (single pass)     & AP (aug-avg $K{=}16$) & $\Delta$ aug \\
\midrule
\tie map             &   16  &      5,824 & 0.0622 $\pm$ 0.0023  & ---                   & ---           \\
\tie oap             &   16  &      5,824 & 0.0605 $\pm$ 0.0022  & ---                   & ---           \\
\win random\_augmented &   16  &      5,824 & \textbf{0.0636} $\pm$ 0.0029 & 0.0637 $\pm$ 0.0029 & $+$0.0001    \\
\midrule
\tie map             &  128  &    203,392 & 0.1705 $\pm$ 0.0141  & ---                   & ---           \\
\tie oap             &  128  &    203,392 & 0.1754 $\pm$ 0.0073  & ---                   & ---           \\
\win random\_augmented &  128  &    203,392 & \textbf{0.1766} $\pm$ 0.0062 & 0.1766 $\pm$ 0.0062 & $-$0.0000    \\
\midrule
\tie map             &  512  &  2,975,872 & 0.2344 $\pm$ 0.0078  & ---                   & ---           \\
\tie oap             &  512  &  2,975,872 & 0.2371 $\pm$ 0.0076  & ---                   & ---           \\
\win random\_augmented &  512  &  2,975,872 & \textbf{0.2395} $\pm$ 0.0139 & 0.2400 $\pm$ 0.0140 & $+$0.0005    \\
\bottomrule
\end{tabular}
\caption{Sweep~2, $k=8$, ogbg-molpcba, mean $\pm$ std over 3 seeds.
\texttt{random\_augmented} wins at \emph{every} hidden dim. The
$h\!=\!512$ verdict from sweep~1 (where \texttt{oap} won) flipped: at
$k\!=\!8$ random\_augmented edges out \texttt{oap} by $+0.0024$ AP.
Test-time averaging again contributes $\leq +0.0005$.}
\label{tab:k8}
\end{table}

\subsection{Side-by-side: $k=3$ vs $k=8$}

\begin{table}[h]
\centering\small
\begin{tabular}{r l l l l}
\toprule
$h$  & \multicolumn{2}{c}{$k=3$}        & \multicolumn{2}{c}{$k=8$}                \\
     & best canon         & winner       & best canon         & winner             \\
\midrule
 16  & oap   $0.0641$     & rand\_aug    & map   $0.0622$     & rand\_aug          \\
128  & map   $0.1704$     & rand\_aug    & oap   $0.1754$     & rand\_aug          \\
\textbf{512} & \textbf{oap   $\mathbf{0.2431}$} & \textbf{oap} & oap   $0.2371$     & \textbf{rand\_aug} \\
\bottomrule
\end{tabular}
\caption{The $h\!=\!512$ verdict flips between sweeps: \texttt{oap} on top
at $k\!=\!3$, \texttt{random\_augmented} on top at $k\!=\!8$. The PE
width matters; widening from $3$ to $8$ eigenvectors changes the
canonicalization-vs-augmentation winner at the largest model.}
\label{tab:compare}
\end{table}

\subsection{Visual summary}

\begin{figure}[h]
\centering
\begin{minipage}{0.48\textwidth}
\centering
\includegraphics[width=\linewidth]{../plots/ap_vs_hdim.png}
\caption*{$k=8$, single-pass eval}
\end{minipage}\hfill
\begin{minipage}{0.48\textwidth}
\centering
\includegraphics[width=\linewidth]{../plots/ap_vs_hdim_augmented.png}
\caption*{$k=8$, with aug-averaging $K=16$}
\end{minipage}
\caption{$k\!=\!8$ test AP vs.\ hidden dim. Left: \texttt{map}, \texttt{oap},
\texttt{random\_augmented} single-pass. Right: same plus the aug-averaged
\texttt{random\_augmented} curve (dashed). Error bars are $\pm 1\,\sigma$
over $3$ seeds. \texttt{random\_augmented} stays on top at every $h$;
the aug-averaging curve sits essentially on top of the single-pass curve.}
\label{fig:plots}
\end{figure}

%==============================================================================
\section{Interpretation}
%==============================================================================

\subsection{Does canonicalization perform better? --- No, with one corner.}

Looking across the $54$ runs, \texttt{random\_augmented} is the best (or
tied) method in $5$ of the $6$ \,$(k,h)$\, cells. The single exception is
$k\!=\!3$ / $h\!=\!512$, where \texttt{oap} edges out
\texttt{random\_augmented} by $+0.0033$ AP --- a real but narrow win.
At $k\!=\!8$ that win disappears.

If you're deciding between canonicalization and augmentation for a new
project on similar data, the safer default is \textbf{random augmentation
with $k\!\geq\!8$}.

\subsection{Why does the $h\!=\!512$ verdict flip?}

Two stories that both fit the data:
\begin{enumerate}\itemsep0pt
\item \emph{Averaging-out story.} At $k\!=\!8$ training has to internalize
      invariance across $256$ sign patterns versus only $8$ at $k\!=\!3$.
      The wider implicit-invariance training produces an encoder that
      generalizes better at test time.
\item \emph{Information story.} $8$ eigenvectors carry materially more
      structural signal. Canonicalization rules trade that off for
      determinism; augmentation lets the network exploit the richer signal
      as long as it's invariance-trained.
\end{enumerate}
The data does not separate them. Both predict that going from $k\!=\!3$ to
$k\!=\!8$ helps random\_augmented at the high-capacity end --- which is
what we see.

\subsection{Why is the test-time averaging $\Delta$ so small?}

By design. A model trained for tens of epochs with random-sign-flip
augmentation becomes \emph{near-invariant} to a single sign flip, so
averaging $K$ predictions adds little new information. The biggest
$\Delta$ we see is $+0.0005$ AP at $k\!=\!8$ / $h\!=\!512$ --- safe and
free, but not the lever that decides anything. \emph{Test-time averaging
matters most when the model under-fits the invariance during training};
once it's trained well it converges to a single-prediction mode.

\subsection{When does \texttt{oap} still win?}

In the narrow $k\!=\!3$ / $h\!=\!512$ corner. \texttt{oap}'s seed variance
there is $\pm 0.0022$ vs.\ random\_augmented's $\pm 0.0117$ --- $5\times$
tighter. With $3$ seeds, that means \texttt{oap}'s confidence interval
is narrower and the modal seed beats random\_augmented's modal seed. With
more seeds and at $k\!=\!8$ the picture inverts. So the $k\!=\!3$/$h\!=\!512$
\texttt{oap} win is best read as ``\texttt{oap} is the lower-variance
choice if you're forced to use $k\!=\!3$,'' not ``canonicalization beats
augmentation in general.''

%==============================================================================
\section{Main conclusions}
%==============================================================================

\begin{enumerate}
\item \textbf{Canonicalization does not generally beat random augmentation
      on ogbg-molpcba.} Across the $6$ tested $(k,h)$ cells,
      \texttt{random\_augmented} wins or ties in $5$.

\item \textbf{Use $k\geq 8$.} At $k\!=\!3$ canonicalization (specifically
      \texttt{oap}) has an edge at the largest hidden dim. At $k\!=\!8$
      that edge disappears.

\item \textbf{Test-time ambiguity-group averaging is safe-and-free, not
      transformative.} It contributes $+0.0001$ to $+0.0005$ AP and never
      hurts. Use it if it's free; do not expect it to flip a verdict.

\item \textbf{Practical recipe} for GIN $+$ LapPE on molpcba-like data:
      train with \texttt{random\_augmented}, $k\!=\!8$, hidden dim $\geq
      128$, patience $15$, and run $K\!=\!8$--$16$ test-time averaging.
\end{enumerate}

\paragraph{Caveats.}
\begin{itemize}\itemsep0pt
\item All numbers are 3-seed averages. Several gaps are within $1$--$2\sigma$.
\item Restricted to GIN, ogbg-molpcba, $k\!\in\!\{3,8\}$,
      $h\!\in\!\{16,128,512\}$. Whether the trend extends to other
      datasets, larger $k$, transformer backbones, or atomic-data tasks is
      not tested here.
\item \texttt{spielman}, \texttt{maxabs}, \texttt{abs}, \texttt{none} were
      out of scope for these two sweeps. See sibling
      \texttt{results/nadav\_improvements/} for those methods under a
      different protocol.
\end{itemize}

\end{document}
